\subsection{不可分解模与原始模}

设 A 是一个环。回忆下面的定义（VII, §4, No.~8, 第 23 页，定义 3）。

定义 2. —— 若 A-模 M 不是一族异于 0 与 M 的子模的直和，则称 A-模 M 为不可分解模。

由 II, §1, No.~8, 第 209 页命题 12 的推论，下列性质等价：

\begin{enumerate}
\def\labelenumi{\alph{enumi})}
\item
  A-模 M 不可分解。
\item
  A-模 M 非零，且 M 的每个作为直因子的子模都等于 0 或 M。
\item
  A-模 M 非零，且环 \(\operatorname{End}_{\mathrm{A}}(\mathrm{M})\) 不含异于 0 和 \(1_{M}\) 的幂等元。
\end{enumerate}

特别地，由于 A-模 \(A_{s}\) 的自同态环同构于 A 的反环，我们可见 A-模 \(A_{s}\) 不可分解当且仅当环 A 非零且其仅有的幂等元是 0 和 1。

例. —— 设环 A 是主理想整环。不可分解的有限生成 A-模是同构于 A 或 \(A/p^{n}A\) 的那些 A-模，其中 p 是 A 的不可约元，n 为整数 \textgreater0（VII, §4, No.~8, 第 24 页，命题 8）。

命题 3. —— 诺特的或阿廷的 A-模 M 是不可分解子模的有限族的直和。

先证明 M 的每个非零子模 P 都有一个不可分解的直因子。若不然，则 P 的每个直因子子模都是可分解的；于是用归纳法，可对每个 \(n \in N\) 构造 P 的非零子模 \(N_{n}^{\prime}\) 与 \(N_{n}^{\prime\prime}\)，使得 \(P = N_{0}^{\prime} \oplus N_{0}^{\prime\prime}\) 且 \(N_{n-1}^{\prime} = N_{n}^{\prime} \oplus N_{n}^{\prime\prime}\)（\(n \geqslant 1\)）。但这样一来，子模序列 \(N_{0}^{\prime\prime} + \cdots + N_{n}^{\prime\prime}\) 将严格递增，而子模序列 \(N_{n}^{\prime}\) 将严格递减。于是模 M 既非诺特的亦非阿廷的，与假设矛盾。

现在设 M 不是不可分解子模的有限族的直和。用归纳法，我们构造 M 的不可分解子模 \(P_{n}^{\prime\prime}\) 与 M 的非零子模 \(P_{n}^{\prime}\)（对每个 \(n \in N\)），使得 \(M = P_{0}^{\prime} \oplus P_{0}^{\prime\prime}\) 且 \(P_{n-1}^{\prime} = P_{n}^{\prime} \oplus P_{n}^{\prime\prime}\)（\(n \geqslant 1\)）。事实上，M 非零，把证明的第一部分应用于 P = M，即得到 \(P_{0}^{\prime}\) 与 \(P_{0}^{\prime\prime}\)。由于模 \(P_{k}^{\prime}\) 与 \(P_{k}^{\prime\prime}\) 对 k \textless{} n 已有定义，由证明的第一部分，存在子模 \(P_{n}^{\prime}\) 与 \(P_{n}^{\prime\prime}\)，使得 \(P_{n-1}^{\prime} = P_{n}^{\prime} \oplus P_{n}^{\prime\prime}\)，其中 \(P_{n}^{\prime\prime}\) 不可分解。因为 M 不是不可分解模的有限族的直和，关系式 \(M = P_{n}^{\prime} \oplus P_{0}^{\prime\prime} \oplus \cdots \oplus P_{n}^{\prime\prime}\) 便蕴含 \(P_{n}^{\prime} \neq 0\)。

子模序列 \(P_{0}^{\prime\prime} \oplus \cdots \oplus P_{n}^{\prime\prime}\) 严格递增，而子模序列 \(P_{n}^{\prime}\) 严格递减。这与 M 是诺特的或阿廷的这一假设相矛盾。

把模分解为不可分解子模的直和时的唯一性问题，将在下一小节中研究。

定义 3. —— 若一个模的自同态环是局部环，则称该模为原始模 \(^{(1)}\)。

按定义，局部环不化为 0；因此，原始模非零。此外，A-模 \(A_{s}\) 是原始模当且仅当环 A 是局部环。

命题 4. —— a) 原始模是不可分解模。

\begin{enumerate}
\def\labelenumi{\alph{enumi})}
\setcounter{enumi}{1}
\tightlist
\item
  有限长度的不可分解模是原始模。
\end{enumerate}

设 M 是一个 A-模。设 M 是原始模；设 e 是局部环 \(\operatorname{End}_{\mathrm{A}}(\mathrm{M})\) 中的一个幂等元。由于 \(e^{2}=e\)，或者 e 可逆从而 e=1，或者 1-e 可逆从而 e=0。这就证明 M 不可分解（VIII，第 30 页）。

现在设 M 不可分解且长度有限。由 VIII，第 27 页命题 2, c)，M 的每个自同态都可逆或幂零；于是环 \(\operatorname{End}_{\mathrm{A}}(\mathrm{M})\) 是局部环（VIII，第 26 页，例 2）。

Z-模 Z 不可分解、是诺特的，但不是阿廷的。它的自同态环同构于 Z，因而不是局部环。于是 Z 不是原始 Z-模。

*设 \(p\) 是素数。\(\mathbf{Z}\) -模 \(\mathbf{Q}_p / \mathbf{Z}_p\) 的自同态环同构于局部环 \(\mathbf{Z}_p\)（参见 VII, §4, 第 65 页，习题 13）；因此它是原始 \(\mathbf{Z}\) -模。

一个内射模不可分解当且仅当它是原始模（X, §1, n° 9, 第 21 页，命题 14）。*

